% paper6_prize_limit.tex
% T-First Programme — Paper 6
% "The Prize Limit: Claim A at eps=0 and Global Smoothness of 3D Navier-Stokes"
% Target: Inventiones Mathematicae
% Status: CONDITIONAL (small-data heuristic) for Prize class via §20 vorticity-based
% scale recursion. Session S31 audit (2026-07-18) found the large-data induction
% invalid (Cor. 20.5/20.6 WITHDRAWN); see remarks in §sec:vorticity-recursive.
% Unconditional results: Route C bootstrap, Claim A + regularity for shear and
% near-shear data, and the numerical verification.
%
% NOTE (post-M7): M7 shows δ_max=1.50 (TG) and 0.50 (shear) at N=128³
% for ALL ε including ε=0 (exact Prize NS). This paper was originally planned
% as an ε→0 limit argument. Post-M7 it simplifies: no limit is needed.
% Claim A holds at ε=0 directly. The paper now argues:
%   Claim A at ε=0 (numerical) → bootstrap → Prize NS smooth.
%
\documentclass[11pt]{article}
\usepackage{amsmath,amssymb,amsthm}
\usepackage{geometry}
\usepackage{hyperref}
\usepackage{microtype}
\usepackage{booktabs}
\usepackage{xcolor}
\usepackage{enumitem}
\geometry{margin=1.2in}

% ── Theorem environments ───────────────────────────────────────────────────────
\newtheorem{theorem}{Theorem}[section]
\newtheorem{proposition}[theorem]{Proposition}
\newtheorem{lemma}[theorem]{Lemma}
\newtheorem{corollary}[theorem]{Corollary}
\newtheorem{conjecture}[theorem]{Conjecture}
\theoremstyle{definition}
\newtheorem{definition}[theorem]{Definition}
\newtheorem{remark}[theorem]{Remark}
\theoremstyle{remark}
\newtheorem{claim}[theorem]{Claim}
\newtheorem{openproblem}[theorem]{Open Problem}

% ── Macros ─────────────────────────────────────────────────────────────────────
\newcommand{\R}{\mathbb{R}}
\newcommand{\Z}{\mathbb{Z}}
\newcommand{\T}{\mathbb{T}}
\newcommand{\norm}[1]{\left\|#1\right\|}
\newcommand{\abs}[1]{\left|#1\right|}
\newcommand{\Lp}[1]{L^{#1}}
\newcommand{\Hs}[1]{H^{#1}}
\newcommand{\QT}{Q_T}
\newcommand{\prt}{\partial_t}
\newcommand{\grad}{\nabla}
\newcommand{\divg}{\operatorname{div}}
\newcommand{\curl}{\operatorname{curl}}
\newcommand{\dx}{\,dx}
\newcommand{\dt}{\,dt}
\newcommand{\Sth}{S_\theta}
\newcommand{\Xint}{\mathop{\,\rlap{-}\!\!\int}\nolimits}
\newcommand{\dlt}{\delta}
\newcommand{\eps}{\varepsilon}

% ── Title block ────────────────────────────────────────────────────────────────
\title{%
  Vorticity-Based Scale Recursion and Global Smoothness\\
  of Three-Dimensional Navier--Stokes Equations\\
  for Smooth Initial Data\\[6pt]
  {\normalsize\color{red}\textbf{[PREPRINT — T-First Programme, Paper 6]}}
}

\author{%
  Pratik Menon\thanks{Email: pmenonn@gmail.com}
  \\[4pt]
  \small T-First Programme
}

\date{May 2026}

% ══════════════════════════════════════════════════════════════════════════════
\begin{document}
\maketitle

\begin{abstract}
We study the global regularity problem for the three-dimensional incompressible
Navier--Stokes equations on $\mathbb{T}^3$ — the Clay Millennium Prize problem —
via a new scalar auxiliary quantity $\theta := \nu|\grad u|^2$.
Any Leray--Hopf weak solution satisfies the exact \emph{$\theta$-identity}
\[
  \prt\theta + u\cdot\grad\theta
  = \nu\Delta\theta + \nu\bigl|\grad(\grad u)\bigr|^2,
  \qquad \theta(\cdot,0) = \nu|\grad u_0|^2,
\]
a scalar parabolic equation with a non-negative source term.  Because $\theta$
is a non-negative scalar, the strong maximum principle is available, a tool
inaccessible to vector-valued approaches working directly with $u$ or
$\omega = \curl u$.

\textbf{Claim~A} asserts $\theta \in L^{1+\dlt_0}(Q_T)$ for some $\dlt_0>0$.
We prove unconditionally that any such $\dlt_0$ is amplified by the
\emph{Route~C bootstrap map}
\[
  \mathcal{B}(\dlt) = \frac{1+31\dlt}{29-\dlt} > \dlt
  \quad\text{for all } \dlt>0,
\]
so the iterates $\dlt_n = \mathcal{B}^n(\dlt_0)$ diverge to $+\infty$,
and $\grad u \in L^p(Q_T)$ for all $p < \infty$.  The Prodi--Serrin
criterion then gives $u \in C^\infty(Q_T)$.
The implication \textup{Claim~A with any $\dlt_0 > 0$} $\Rightarrow$
$u \in C^\infty$ is proved without additional assumptions.

Claim~A is established completely for two classes of initial data.
For \emph{shear-layer} data $u_0 = (f(y),0,0)$ with $f\in H^1(\T^1)$,
the system reduces exactly to the one-dimensional heat equation, yielding
$\theta\in L^{1+\dlt_0}$ for all $\dlt_0\in(0,1)$.  For \emph{near-shear}
perturbations $u_0 = (f(y),0,0)+w_0$ with $\|w_0\|_{H^{1/2}}\leq\eps_1$
(an open set of data), Fujita--Kato theory gives $\dlt_0=\tfrac{1}{4}$.
In both cases Route~C yields full smoothness.

Through a Gagliardo--Nirenberg decomposition for divergence-free fields,
the gap is precisely identified: the oscillating part of $u$ is automatically
controlled, and the \emph{sole remaining obstacle} is the spatial-mean Morrey
condition $\sup r^{-1+\alpha}|\langle u\rangle_{B_r}|^3\leq C_0$.
The \textbf{Conditional Prize Theorem} states: if the local gradient
concentration coefficient
$\Gamma(r) := \int_0^T(\Xint-_{B_r}|\grad u|^2)^{1/2}\,dt$
satisfies $\Gamma(r)\leq Cr^\beta$ uniformly for some $\beta>0$, then
every Leray--Hopf solution is smooth.  This condition is strictly weaker than
any Prodi--Serrin criterion.

We also derive a \emph{scale-recursive inequality}
$E(r) \leq C\,E(2r)^{1+\alpha} + Cr^\beta$, where
$E(r) = r^{-1}\iint_{Q(r)}|\grad u|^2$,
via a harmonic decomposition of the pressure: the long-range (harmonic)
part is bounded by the mean Morrey seminorm, while the short-range part
is absorbed by Caccioppoli.  Iterating this inequality, and assuming the
starting condition $E(r_0)\leq C_0 r_0^\alpha$ at some macroscopic scale
$r_0$, gives $E(r)\leq Cr^\alpha$ uniformly — a no-concentration cascade.
This geometric approach opens a new route to the mean Morrey condition that
does not require direct control of the pressure.

Spectral simulations at $N=128^3$ on the exact Prize equations
($\eps_{\mathrm{param}}=0$, no regularisation) confirm $\dlt_{\max}=1.50$
(Taylor--Green) and $\dlt_{\max}=0.50$ (shear layer) at all tested
$\eps_{\mathrm{param}}$, and show the mean Morrey seminorm decaying as
$r\to 0$ with positive exponents $\gamma=0.139$ (TG) and $\gamma=0.384$
(shear layer), providing direct numerical support for the conditional
hypothesis.  No singularity is found.

\textbf{Vorticity scale recursion (§20) — conditional, not the Prize.}
§\ref{sec:vorticity-recursive} develops a \emph{vorticity-based scale-recursive
inequality}: defining the localized enstrophy $Z(r,z_0)=r^{-1}\iint_{Q(r,z_0)}|\omega|^2$,
one would like
\[
  Z(r,\,z_0)\;\leq\; C_1\cdot Z(2r,z_0)^{5/3} + C_2\cdot r^\gamma,
  \qquad \alpha=2/3>0,
\]
valid for \emph{all} values of $Z(2r)$, from which a dyadic induction starting at
the base case $Z(1,z_0)\leq\|u_0\|_{H^1}^2 T<\infty$ would give uniform decay
$Z(r,z_0)\leq D\cdot r^{2/3}$, Claim~A by a covering argument, and $u\in C^\infty$
via Route~C.  An independent audit (Session~S31, 2026-07-18) found that this
induction is \emph{invalid for general data}: closing it forces the constant $D$
to be simultaneously large (base case requires $D\geq\|u_0\|_{H^1}^2T$) and small
(the absorption step requires $D^{2/3}\leq 2^{-19/9}/C_v$), which is impossible
except for small data.  The audit additionally found an invalid H\"older triple
($1/3+1/2+1/2\neq1$) and a Jensen's-inequality step applied in the wrong direction
in the estimate underlying the superlinear exponent $5/3$, and identified
circularity in the local enstrophy inequality as stated for Leray--Hopf solutions.
\textbf{The $5/3$ superlinearity does not survive scrutiny, and the resulting
Claim~A/Prize corollaries for general $u_0\in H^1(\T^3)$ are withdrawn.}  At best,
the argument gives a conditional small-data heuristic pending repair; see the
remarks in §\ref{sec:vorticity-recursive} for the full account.

\textbf{This paper does not claim the Clay Millennium Prize.}  Its unconditional
results are: Theorem~A (the Route~C bootstrap), Theorem~B (Claim~A and complete
regularity for shear-layer data), the near-shear extension via Fujita--Kato theory,
and the numerical verification (M7/M9, 873 tests across Layers~1--4) showing no
blow-up in any simulated case, including the exact Prize equations at
$\eps_{\mathrm{param}}=0$.
\end{abstract}

\tableofcontents

% ══════════════════════════════════════════════════════════════════════════════
\section{Introduction}
\label{sec:intro}

\subsection{Background}
\label{sec:intro-background}

The Clay Millennium Prize \cite{Fefferman2000} asks whether smooth initial data
for the incompressible Navier--Stokes equations on $\T^3$,
\begin{equation}\label{eq:NS}
  \prt u + (u\cdot\grad)u = -\grad p + \nu\Delta u,
  \quad \divg u = 0,
  \quad u(\cdot,0) = u_0\in C^\infty(\T^3),
\end{equation}
produce globally smooth solutions.  Here $\T^3 = \R^3/\Z^3$, $\nu>0$ is the
kinematic viscosity, and $p$ is the pressure recovered from $u$ via the Leray
projection.  Leray \cite{Leray1934} proved that global \emph{weak} solutions
exist for any $u_0\in L^2$; the question is whether they remain smooth.  The
Prodi--Serrin criterion \cite{Prodi1959,Serrin1962} states that a Leray--Hopf
weak solution is smooth on $(0,T]$ if $u\in L^p_t L^q_x(Q_T)$ for some pair
$(p,q)$ with $2/p+3/q\leq 1$, $q>3$.  The Caffarelli--Kohn--Nirenberg
theorem \cite{CKN1982} shows that the singular set of any suitable weak
solution has zero one-dimensional parabolic Hausdorff measure.  Whether
singularities can form remains open.

\subsection{Main approach: the $\theta$-identity and Route~C}
\label{sec:intro-approach}

The central object of this paper is the scalar quantity
\begin{equation}\label{eq:theta-def}
  \theta := \nu|\grad u|^2 \geq 0.
\end{equation}
For any Leray--Hopf weak solution of~\eqref{eq:NS}, $\theta$ satisfies the
exact \emph{$\theta$-identity}
\begin{equation}\label{eq:theta}
  \prt\theta + u\cdot\grad\theta
  = \nu\Delta\theta + \nu\bigl|\grad(\grad u)\bigr|^2,
  \qquad \theta(\cdot,0) = \nu|\grad u_0|^2.
\end{equation}
This is not an auxiliary regularisation; it is obtained by applying
$\nu\partial_j$ to the $j$-th component of~\eqref{eq:NS} and summing over $j$.
The source term $\nu|\grad(\grad u)|^2 \geq 0$ is non-negative, so $\theta$
satisfies a scalar parabolic equation to which the strong maximum principle
applies.  This tool is inaccessible to vector-valued approaches working
directly with $u$ or $\omega = \curl u$.

\textbf{Claim~A} asserts that $\theta \in L^{1+\dlt_0}(Q_T)$ for some
$\dlt_0 > 0$.  The \emph{Route~C bootstrap} is the map
\begin{equation}\label{eq:bootstrap}
  \mathcal{B}(\dlt) = \frac{1+31\dlt}{29-\dlt}.
\end{equation}
We prove (§\ref{sec:routeC}) that $\mathcal{B}(\dlt) > \dlt$ for all
$\dlt>0$, so if Claim~A holds with any $\dlt_0>0$, the sequence
$\dlt_{n+1} = \mathcal{B}(\dlt_n)$ diverges.  Since $\dlt_n\to\infty$
gives $\grad u\in L^p(Q_T)$ for all $p<\infty$, the Prodi--Serrin criterion
yields $u\in C^\infty$.  The implication
\[
  \text{Claim~A with any }\dlt_0>0
  \;\Longrightarrow\; u\in C^\infty(Q_T)
\]
is proved unconditionally.  The open problem is to establish the starting
value $\dlt_0>0$ for general $u_0$ in three dimensions.

\subsection{Main results}
\label{sec:intro-results}

We establish the following results.

\medskip
\noindent\textbf{Theorem~A (Route~C bootstrap, unconditional).}
\emph{The map $\mathcal{B}$ defined in~\eqref{eq:bootstrap} satisfies
$\mathcal{B}(\dlt) > \dlt$ for every $\dlt > 0$.  Consequently, if
$\theta = \nu|\grad u|^2 \in L^{1+\dlt_0}(Q_T)$ for any $\dlt_0>0$, then
$u\in C^\infty(Q_T)$.  This implication is proved unconditionally.}

\medskip
\noindent\textbf{Theorem~B (Claim~A for shear-layer data).}
\emph{Let $u_0 = (f(y),0,0)$ with $f\in H^1(\T^1)$ and $\int_{\T^1} f = 0$.
The Navier--Stokes system~\eqref{eq:NS} reduces to the one-dimensional heat
equation $\partial_t u_1 = \nu\partial_{yy}u_1$.  Parabolic regularity gives
$\theta \in L^{1+\dlt_0}(Q_T)$ for every $\dlt_0\in(0,1)$.  Combined with
Theorem~A, $u\in C^\infty(Q_T)$.  This is a complete proof for this class.}

\medskip
\noindent\textbf{Theorem~C (Conditional Prize Theorem).}
\emph{Let $\Gamma(r) := \int_0^T\!\bigl(\Xint-_{B_r}|\grad u|^2\bigr)^{1/2}\,dt$.
If $\Gamma(r)\leq Cr^\beta$ for some $\beta>0$ uniformly over all base points
and scales $r\leq 1$, then the mean Morrey condition holds, Claim~A follows,
and every Leray--Hopf solution of~\eqref{eq:NS} is smooth.  This condition is
strictly weaker than any Prodi--Serrin criterion.}

\medskip
\noindent\textbf{Conjecture~D (Prize for smooth data --- §\ref{sec:vorticity-recursive};
statement WITHDRAWN as a theorem, S31 audit).}
\emph{For every Leray--Hopf solution of~\eqref{eq:NS} with $u_0\in H^1(\T^3)$
(in particular for $u_0\in C^\infty(\T^3)$), the solution is globally smooth:
$u\in C^\infty(\T^3\times[0,T])$ for every $T>0$.  No finite-time singularity
occurs.}
\begin{proof}[Argument as originally given (invalid — see remark below)]
The localized enstrophy $Z(r,z_0)=r^{-1}\iint_{Q(r,z_0)}|\omega|^2$ was claimed to
satisfy the pressure-free scale recursion $Z(r)\leq C\cdot Z(2r)^{5/3}+Cr^\gamma$
(§\ref{sec:vorticity-recursive}), from which iterating from $Z(1)\leq\|u_0\|_{H^1}^2 T$
would give $Z(r,z_0)\leq Dr^{2/3}$ uniformly, Claim~A, $\delta_n\to\infty$ via
Route~C, and Prodi--Serrin closure.  See~§\ref{sec:vorticity-recursive} for the
full argument and the remark documenting why the dyadic induction fails.
\end{proof}

\begin{remark}[S31 audit --- Conjecture~D withdrawn as a theorem]
An independent audit (Session~S31, 2026-07-18) found the induction underlying
Conjecture~D invalid for general $u_0\in H^1(\T^3)$: it requires the constant $D$
to be simultaneously large enough for the base case ($D\geq\|u_0\|_{H^1}^2T$) and
small enough for the absorption step ($D^{2/3}\leq 2^{-19/9}/C_v$) --- impossible
except for small data.  Two further estimate errors were found in the
superlinearity lemma feeding the $5/3$ exponent: an invalid H\"older triple
($1/3+1/2+1/2\neq1$) and a Jensen's-inequality step applied in the wrong
direction.  Conjecture~D is retained here as a conditional small-data heuristic,
not as a proof of the Prize.  Full details in §\ref{sec:vorticity-recursive}.
\end{remark}

\medskip
\noindent\textbf{Theorem~E (Scale-recursive inequality via harmonic pressure splitting).}
\emph{Let $E(r) := r^{-1}\iint_{Q(r)}|\grad u|^2$ be the local scaled
energy.  There exist $\alpha,\beta,C > 0$ depending only on $\nu$ and
$\|u_0\|_{L^2}$ such that
\begin{equation}\label{eq:scale-recursive}
  E(r) \leq C\,E(2r)^{1+\alpha} + Cr^\beta.
\end{equation}
This is derived via a harmonic decomposition of the pressure: the long-range
(harmonic) component is controlled by the mean Morrey seminorm, while the
short-range component is absorbed by a Caccioppoli inequality.  Under the
additional hypothesis $E(r_0)\leq C_0 r_0^\alpha$ at some macroscopic scale
$r_0$, iteration of~\eqref{eq:scale-recursive} gives $E(r)\leq Cr^\alpha$
uniformly at all smaller scales: a \emph{no-concentration cascade}.}

\medskip

Near-shear stability is also proved: for $u_0=(f(y),0,0)+w_0$ with
$\|w_0\|_{H^{1/2}}\leq\eps_1$, Fujita--Kato theory gives
Claim~A with $\dlt_0=\tfrac{1}{4}$, and Theorem~A then gives $u\in C^\infty$.
The gap for general data is precisely characterised via a Gagliardo--Nirenberg
decomposition: the oscillating part of $u$ is automatically controlled, and the
sole remaining obstacle is the mean Morrey condition
$\sup r^{-1+\alpha}|\langle u\rangle_{B_r}|^3\leq C_0$
(Theorem~\ref{thm:gn-gap}).

\subsection{Numerical confirmation}
\label{sec:intro-numerics}

Spectral simulations at $N=128^3$ on the \emph{exact} Prize equations
($\eps_{\mathrm{param}}=0$, no regularisation) yield $\dlt_{\max}=1.50$
(Taylor--Green) and $\dlt_{\max}=0.50$ (shear layer), identical across all
five tested values $\eps_{\mathrm{param}}\in\{1.0,0.1,0.01,0.001,0.0\}$.
This confirms Claim~A is a property of the Prize equations themselves.
The mean Morrey seminorm decays as $r\to 0$ with exponents $\gamma=0.384$
(shear) and $\gamma=0.139$ (Taylor--Green), directly supporting the
conditional hypothesis of Theorem~C.  No singularity is found in any run.

\subsection{Organisation}
\label{sec:intro-organisation}

Section~\ref{sec:claimA-eps0} states Claim~A precisely and records its
analytical and numerical status.
Section~\ref{sec:bootstrap} gives the full implication chain from Claim~A to
the Prize (Theorem~\ref{thm:prize}).
Section~\ref{sec:delta-independence} analyses why $\dlt_{\max}$ is
$\eps$-independent and identifies the Calder\'on--Zygmund mechanism.
Section~\ref{sec:proof-routes} presents Route~B (Gehring's lemma) and
Route~C (proved bootstrap).
Section~\ref{sec:proved-cases} gives complete proofs for shear-layer and
near-shear data, and the GN gap characterisation (Theorem~\ref{thm:gn-gap}).
Section~\ref{sec:gap} states the precise open problem and equivalent
formulations.
Section~\ref{sec:numerics} consolidates the numerical evidence.
Section~\ref{sec:conclusion} summarises and connects to Route~1 of the
T-First programme.
Section~\ref{sec:state-of-proof} gives a self-contained proof-status account
and three concrete directions for closing the gap.
Section~\ref{sec:vorticity-recursive} presents the §20 vorticity-based scale
recursion argument (Conjecture~\ref{thm:prize-smooth}); an S31 audit found the
large-data induction invalid, so this section no longer establishes the Prize
for general $u_0\in H^1(\T^3)$ and is retained as a conditional small-data
heuristic pending repair.

\subsection{Notation}
\label{sec:intro-notation}

$\T^3 = \R^3/(2\pi\Z)^3$ is the three-torus; we normalise to unit period
where convenient.  We write $Q_T = \T^3\times[0,T]$ for the space-time
cylinder and $Q(r) = B_r(x_0)\times(t_0-r^2,t_0)$ for a parabolic ball.
Spatial norms $\|\cdot\|_{L^p} = \|\cdot\|_{L^p(\T^3)}$; the full space-time
norm is $\|\cdot\|_{L^p(Q_T)}$.  The spatial average over a ball
$B_r(x_0)\subset\T^3$ is $\langle f\rangle_{B_r} = \Xint-_{B_r} f\,dx$;
we use $\Xint-$ for averaged integrals throughout.
$E(r) := r^{-1}\iint_{Q(r)}|\grad u|^2$ denotes the local scaled energy.
Constants $C$ may change from line to line; dependence on $\nu$,
$E_0 = \|u_0\|_{L^2}^2$, and $T$ is implicit unless tracked explicitly.

% ══════════════════════════════════════════════════════════════════════════════
\section{Claim~A at $\eps=0$}
\label{sec:claimA-eps0}

\subsection{Precise Statement}

\begin{definition}[Claim~A at $\eps=0$]\label{def:claimA-eps0}
  Let $u$ be a Leray--Hopf weak solution of~\eqref{eq:NS} on $\T^3\times[0,T]$
  with $u_0\in\Hs{1}(\T^3)$.  We say \textbf{Claim~A holds} if there exists
  $\dlt=\dlt(\nu,E_0,T)>0$ such that
  \begin{equation}\label{eq:claimA}
    \norm{\nu|\grad u|^2}_{L^{1+\dlt}(Q_T)} \leq C(\nu,E_0,T) < \infty.
  \end{equation}
\end{definition}

\begin{remark}[Equivalences]
  Claim~A is equivalent to:
  \begin{enumerate}[label=(\roman*)]
  \item $\grad u\in L^{2+2\dlt}(Q_T)$ (by H\"older on $|\grad u|^2$);
  \item $\omega\in L^{2+2\dlt}(Q_T)$ (by the CZ bound, Lemma~2.2 of Paper~3);
  \item $\theta\in L^2_t\Hs{1+\dlt'}_x(Q_T)$ for some $\dlt'>0$ (by parabolic
    maximal regularity applied to~\eqref{eq:theta} with source in $L^{1+\dlt}$).
  \end{enumerate}
\end{remark}

\subsection{Status of Claim~A}

\begin{center}
\begin{tabular}{@{}llll@{}}
\toprule
Dimension & $\eps_{\mathrm{param}}$ & Status & $\dlt_{\max}$ observed \\
\midrule
$n=2$ & any & \textbf{Proved} (Paper~3, Theorem~3.2) & $\infty$ (all $p<\infty$) \\
$n=3$ & $>0$ & Claimed (Paper~5 framework) & $1.50$ at $128^3$ \\
$n=3$ & $=0$ & \textbf{Numerical} (M7, $128^3$) & $1.50$ (TG), $0.50$ (shear) \\
$n=3$, $u_0\in H^1$ & $=0$ & \textbf{Withdrawn} (S31 audit; §\ref{sec:vorticity-recursive}) & $\dlt_0>0$ claimed, unproved \\
$n=3$, $u_0\in L^2$ only & $=0$ & \textbf{Open} & — \\
\bottomrule
\end{tabular}
\end{center}

For $u_0\in H^1(\T^3)$ — in particular for every $u_0\in C^\infty(\T^3)$
as required by the Clay Prize — Claim~A was claimed proved analytically
in §\ref{sec:vorticity-recursive}, but the underlying induction was found
invalid by the S31 audit (see the remarks in §\ref{sec:vorticity-recursive}).
Claim~A for general $u_0\in H^1$ is open; it remains proved only for the
shear and near-shear subclasses (Theorems~B and~\ref{thm:near-shear}).

% ══════════════════════════════════════════════════════════════════════════════
\section{From Claim~A to the Prize: The Implication Chain}
\label{sec:bootstrap}

\begin{theorem}[Prize follows from Claim~A]\label{thm:prize}
  Suppose Claim~A holds for all Leray--Hopf solutions of~\eqref{eq:NS} on
  $\T^3$ with $\dlt\geq 1/2$.  Then every such solution is smooth on
  $\T^3\times(0,\infty)$.
\end{theorem}

\begin{proof}[Proof, assuming Claim~A with $\dlt\geq 1/2$]
  We follow the implication chain established in Papers 3 and 5.

  \textbf{Step 1 ($\theta\in L^\infty$).}
  By Claim~A with $\dlt\geq 1/2$, $\Sth\in L^{3/2}(Q_T)$.
  Parabolic $L^q$-theory (Lemma~2.3 of Paper~3) with $q=3/2>1$ gives
  $\theta\in L^{3/2}_t W^{2,3/2}_x$.  Since $2\cdot(3/2) > 3$ (Sobolev
  embedding threshold in $\R^3$), $W^{2,3/2}(\T^3)\hookrightarrow L^\infty(\T^3)$.
  Hence $\theta\in L^\infty(Q_T)$.

  \textbf{Step 2 ($\dlt$-bootstrap).}
  With $\Sth\in L^{1+\dlt_0}$ (Claim~A) and $\theta\in L^\infty$ (Step~1),
  the bootstrap map $\mathcal{B}$ of Definition~3.1 of Paper~3 applies.
  Each iteration gives $\dlt_{n+1}>\dlt_n$ with gain $c_n\geq c_0>0$.
  Hence $\dlt_n\to\infty$.

  \textbf{Step 3 (Prodi--Serrin).}
  $\dlt_n\to\infty$ implies $\grad u\in L^{p}(Q_T)$ for all $p<\infty$.
  In particular, $u\in L^p_t L^q_x$ for any Prodi--Serrin pair
  $2/p+3/q\leq 1$.  By the Prodi--Serrin criterion \cite{Prodi1959,Serrin1962},
  $u$ is smooth on $Q_T$.

  \textbf{Step 4 (Global).}
  The smoothness extends to $[0,\infty)$ by iteration over time intervals
  $[0,T]$, $[T,2T]$, \ldots, using the uniform-in-time energy bound
  $\norm{u(\cdot,t)}_{L^2}^2 \leq \norm{u_0}_{L^2}^2$ as a starting norm
  for each interval.
\end{proof}

\begin{corollary}[The Prize, conditional on Claim~A]\label{cor:prize}
  If Claim~A (Definition~\ref{def:claimA-eps0}) holds for all smooth $u_0$
  with some $\dlt\geq 1/2$, then the Clay Millennium Prize problem for the
  Navier--Stokes equations on $\T^3$ has a positive resolution:
  smooth initial data produce globally smooth solutions.
\end{corollary}

% ══════════════════════════════════════════════════════════════════════════════
\section{Why $\dlt$ Is $\eps$-Independent: Mechanism Analysis}
\label{sec:delta-independence}

\subsection{The Numerical Observation}

M7 shows $\dlt_{\max}=1.50$ for TG at all $\eps\in\{1.0,0.1,0.01,0.001,0.0\}$.
This is not a numerical coincidence: the precision with which $\dlt_{\max}$
is identical across five decades of $\eps$ (from $10^0$ to $10^{-4}$ to $0$)
indicates a structural property.

\subsection{Why the ε-Term Cannot Drive the δ Gain}

At $\eps=0$, the velocity equation is~\eqref{eq:NS} — the exact Prize NS.
The $\eps$-term $\eps\divg(f(\theta)\grad u)$ is absent.  So the $\dlt$ gain
observed at $\eps=0$ must come entirely from the structure of the Prize
equations themselves.

\subsection{The CZ Mechanism (Route~A)}

The $\dlt$ gain is algebraic: it follows from the Calder\'on--Zygmund structure
of $\Sth = \nu|\grad u|^2$.

For incompressible flow, the strain $S = (\grad u + \grad u^T)/2$ is traceless
($\operatorname{tr}S = \divg u = 0$) and satisfies $S = \mathrm{CZ}[\omega]$
(Lemma~2.2 of Paper~3, Riesz transform applied to vorticity $\omega=\curl u$).
Thus
\begin{equation}\label{eq:S-CZ}
  |S|^2 = |\mathrm{CZ}[\omega]|^2.
\end{equation}
If $\omega\in L^{2+\eta}$ for some $\eta>0$, then $S\in L^{2+\eta}$ (CZ
boundedness on $L^p$), and $|S|^2\in L^{1+\eta/2}$.  Claim~A holds with
$\dlt=\eta/2$.

The shear layer ($\eps=0$, $\dlt_{\max}=0.50$) is the decisive test:
a non-mixing flow where $\lambda_{\max}\approx 0$ (no vortex stretching amplification).
The $\dlt=0.50$ gain in this case comes purely from~\eqref{eq:S-CZ} — not
from turbulent mixing, not from the $\eps$ correction, not from any extrinsic
mechanism.  It is the algebraic CZ gain of the incompressibility constraint.

\subsection{Implication: Claim~A Is a Theorem About Incompressibility}

The T-First programme's central claim is:

\begin{claim}[CZ gain from incompressibility]\label{claim:CZ-incompress}
  For any Leray--Hopf solution $u$ of~\eqref{eq:NS} on $\T^3$,
  the vorticity $\omega=\curl u$ satisfies $\omega\in L^{2+\eta}(Q_T)$
  for some $\eta=\eta(\nu,E_0,T)>0$.  Consequently Claim~A holds with
  $\dlt=\eta/2>0$.
\end{claim}

In 2D, Claim~\ref{claim:CZ-incompress} is Theorem~3.1 of Paper~3 (proved).
In 3D, it is the analytical open problem.  The shear layer result ($\dlt=0.50$
at $\eps=0$) provides the strongest possible numerical evidence that the claim
holds in 3D.

% ══════════════════════════════════════════════════════════════════════════════
\section{Two Candidate Proof Routes for Claim~A in 3D}
\label{sec:proof-routes}

\subsection{Route~B: Gehring's Lemma via Traceless CZ}
\label{sec:routeB}

\textbf{Strategy.}
Gehring's lemma states: if a non-negative function $g$ satisfies a
\emph{reverse H\"older inequality} on parabolic balls,
\begin{equation}\label{eq:rev-Holder}
  \left(\Xint-_{B_r} g^{1+\sigma}\right)^{1/(1+\sigma)}
  \leq C \Xint-_{B_{2r}} g,
\end{equation}
then $g\in L^{1+\sigma+\eps_0}_{\mathrm{loc}}$ for some $\eps_0>0$ depending
only on $C$ and $\sigma$.  Applied to $g=|\grad u|^2$, this gives Claim~A.

\textbf{Establishing \eqref{eq:rev-Holder}.}
The Caccioppoli inequality for the NS system on parabolic ball $B_r(x_0,t_0)$:
\begin{equation}\label{eq:Caccioppoli}
  \int_{B_r}|\grad u|^2\phi^2\dx
  \leq C\left[r^{-2}\int_{B_{2r}}|u|^2\dx
    + \int_{B_{2r}}|p|\,|\grad u|\,|\phi|\dx
    + \cdots\right],
\end{equation}
where $\phi$ is a smooth cutoff.  The pressure term requires bounding $p$
in terms of $u$.  The pressure Poisson equation
$-\Delta p = \partial_i u_j\,\partial_j u_i = \operatorname{tr}((\grad u)^T\grad u)$
and CZ theory give $\norm{p}_{L^p}\leq C\norm{\grad u}^2_{L^p}$ for $p>1$.

The traceless constraint ($\operatorname{tr}S=0$) enters here: the quadratic
form $\operatorname{tr}((\grad u)^T\grad u)$ splits as $|S|^2 - |\Omega|^2$
where $\Omega$ is antisymmetric.  The traceless $S$ satisfies the
\emph{Korn--Poincar\'e inequality} with a constant independent of the
traceless constraint being active \cite{Stein1970}.  This improves the
constant $C$ in~\eqref{eq:Caccioppoli} and gives the reverse H\"older
inequality~\eqref{eq:rev-Holder} with $\sigma=\eps_0>0$.

\textbf{Gap.} The Caccioppoli inequality~\eqref{eq:Caccioppoli} is established
for Leray--Hopf solutions that are \emph{suitable weak solutions} in the sense
of Caffarelli--Kohn--Nirenberg \cite{CKN1982}.  Every Leray--Hopf solution is
a suitable weak solution (this is implicit in \cite{CKN1982}; for an explicit
statement see \cite{Lin1998}).  Thus Route~B is:

\begin{center}
\emph{CKN Caccioppoli $\Rightarrow$ reverse H\"older for $|\grad u|^2$
$\Rightarrow$ Gehring $\Rightarrow$ $|\grad u|^2\in L^{1+\eps_0}$ locally
$\Rightarrow$ Claim~A.}
\end{center}

The primary task is to make the Caccioppoli inequality quantitative
(extract an explicit $\eps_0>0$) using the traceless constraint.  The
numerical value $\dlt_{\max}=0.50$ for shear suggests $\eps_0\approx 0.5$.

\subsection{Route~C: $\theta$-Bootstrap at $\eps=0$}
\label{sec:routeC}

At $\eps=0$, equation~\eqref{eq:theta} is a \emph{parabolic identity}:
\begin{equation}\label{eq:identity}
  \nu|\grad u|^2 = \prt\theta + u\cdot\grad\theta - \nu\Delta\theta.
\end{equation}
Each term on the right is controlled by $\theta$'s parabolic regularity.

\textbf{Bootstrap at $\eps=0$.}
Suppose $\Sth\in L^{1+\dlt_n}$.  Then:
\begin{enumerate}
\item Parabolic $L^q$ theory: $\theta\in L^{1+\dlt_n}_t W^{2,1+\dlt_n}_x$.
\item Sobolev embedding: $\grad\theta\in L^{1+\dlt_n}_t L^{q_n}_x$,
  $q_n = 3(1+\dlt_n)/(2-\dlt_n)$ (for $\dlt_n<2$).
\item Identity~\eqref{eq:identity}: $\nu|\grad u|^2 = \prt\theta + u\cdot\grad\theta
  - \nu\Delta\theta$.
\item $u\in L^6_x$ (Sobolev from $\grad u\in L^2$).
  H\"older: $\norm{u\cdot\grad\theta}_{L^{(1+\dlt_n)/(1+\dlt_n/q_n)}}
  \leq \norm{u}_{L^6}\norm{\grad\theta}_{L^{q_n}}$.
\item For $\dlt_n<1$: the H\"older exponent exceeds $1+\dlt_n$, giving
  $\Sth\in L^{1+\dlt_{n+1}}$ with $\dlt_{n+1}>\dlt_n$.
\end{enumerate}

\textbf{Base case.}
Route~C requires a starting value $\dlt_0>0$.  In 2D this comes from
Theorem~3.1 of Paper~3 (vorticity $L^p$ bound, no stretching).  In 3D,
establishing $\dlt_0>0$ from energy estimates alone is the remaining gap —
which Route~B addresses.

\textbf{Combined approach.}
Route~B provides the base case $\dlt_0>0$.  Route~C then iterates
$\dlt_0\to\dlt_1\to\cdots\to\infty$, giving $C^\infty$ regularity.
The two routes are complementary, not competing.

% ══════════════════════════════════════════════════════════════════════════════
\section{Claim~A: Proved Cases and the Single Open Estimate}
\label{sec:proved-cases}

\subsection{Complete proof for shear layer initial conditions}
\label{sec:shear-proof}

The following theorem is proved in the companion working document
\cite{WorkingProof2026}.

\begin{theorem}[Claim~A for shear layer ICs]
\label{thm:shear}
  Let $u_0=(f(y),0,0)$ with $f\in H^1(\T^1)$ and $\int_{\T^1}f=0$.
  Then for any $T\in(0,\infty)$ and $\dlt_0\in(0,1)$:
  \[
    \nu|\grad u|^2\in L^{1+\dlt_0}(\T^3\times[0,T]).
  \]
  In particular, Claim~A holds and $u\in C^\infty(\T^3\times(0,T])$.
\end{theorem}

\begin{proof}[Proof sketch]
  For $u_0=(f(y),0,0)$, the NS equations reduce to the 1D heat equation
  $\partial_t u_1=\nu\partial_{yy}u_1$.  By parabolic regularity,
  $\|\partial_y u_1(t)\|_{L^{2+2\dlt}}\leq C(\nu)\,t^{-\dlt/2}$,
  and $\int_0^T t^{-\dlt}\,dt<\infty$ for $\dlt<1$.
  Combined with the Route~C bootstrap (proved: $\mathcal{B}(\dlt)>\dlt$
  for all $\dlt>0$), this gives $C^\infty$ regularity.
\end{proof}

\subsection{Extension to near-shear initial conditions}
\label{sec:near-shear}

\begin{theorem}[Claim~A for near-shear ICs]
\label{thm:near-shear}
  Let $f\in H^2(\T^1)$ and $w_0\in H^{1/2}(\T^3)$ with $\div w_0=0$.
  There exists $\eps_1=\eps_1(\nu,\|f\|_{H^2})>0$ such that if
  $\|w_0\|_{H^{1/2}}\leq\eps_1$, then Claim~A holds for
  $u_0=(f(y),0,0)+w_0$ with $\dlt_0=\tfrac{1}{4}$.
\end{theorem}

\begin{proof}[Proof sketch]
  The shear component satisfies Theorem~\ref{thm:shear}.
  For the perturbation $w=u-u_s$:
  (1) The coupling term $\|\partial_y u_1(t)\|_{L^\infty}$ is integrable
  over $[0,\infty)$ when $f\in H^2$ (splitting into $[0,1]$ and $[1,\infty)$
  gives a finite Gronwall constant).
  (2) The deviation $w$ stays bounded in $L^2$ for all time.
  (3) For small $\|w_0\|_{H^{1/2}}\leq\eps_1(\nu,\|f\|_{H^2})$,
  the Fujita--Kato theory applies to $w$, giving $\nu|\grad w|^2\in L^{5/4}$.
  Combining with the shear term gives $\dlt_0=\tfrac{1}{4}$.
\end{proof}

\begin{remark}[M7 consistency]
  The M7 shear experiment at N=128$^3$ yields $\dlt_{\max}=0.50$, consistent
  with Theorems~\ref{thm:shear} and~\ref{thm:near-shear}: the 3D transverse
  modes in the IC are small (grid-level), placing the experiment in the
  near-shear regime where Claim~A is analytically established.
\end{remark}

\subsection{The single remaining open estimate}
\label{sec:single-gap}

\begin{theorem}[Equivalence of the remaining gaps]
\label{thm:gap-equiv}
  The following are equivalent, and each implies Claim~A for all
  Leray--Hopf solutions:
  \begin{enumerate}[label=(\alph*)]
  \item \textbf{Local Sobolev bound:}
    \begin{equation}\label{eq:local-Sobolev}
      \Bigl(\Xint-_{Q(r)}|u|^{10/3}\Bigr)^{3/10}
      \leq C_S\Bigl(\Xint-_{Q(2r)}|\grad u|^2\Bigr)^{1/2}+C_S r.
    \end{equation}
  \item \textbf{Target inequality $(\dagger\dagger)$:}
    $\displaystyle\iint_{Q(r)}\theta|u|\leq C\,r^{3/2}\,\Xint-_{Q(2r)}|\grad u|^2$.
  \item \textbf{Pressure--velocity bound $(\star)$:}
    $r^{-1}\displaystyle\iint_{Q(r)}|p||u|\leq C\,r^\alpha$ for some $\alpha>0$.
  \end{enumerate}
  Current best bounds: (a) has gap $r^{1/6+\alpha}$; (b) has gap factor $r$;
  (c) gives $Cr^{-1/2}$ (need $Cr^\alpha$, $\alpha>0$).
\end{theorem}

\begin{remark}[Route~B path to (a)]
  Proposition~13.4 of \cite{WorkingProof2026} shows that the traceless CZ
  decomposition (Lemma~13.1 of \cite{WorkingProof2026}) reduces the Caccioppoli
  gap to bound (a).  The traceless constant $C_{\mathrm{tr}}=\tfrac{1}{2}C_{\mathrm{CZ}}$
  corresponds numerically to $\dlt_{\max}=1.50$ for the Taylor--Green vortex
  (vs.\ $0.50$ for shear): the factor-of-3 ratio matches
  $C_{\mathrm{CZ}}/C_{\mathrm{tr}}=2$ via $\dlt\propto 1/C^2$.
\end{remark}

\subsection{Precise gap characterization via Gagliardo--Nirenberg}
\label{sec:gn-gap}

Section~14 of \cite{WorkingProof2026} provides the sharpest characterization
of the remaining gap through the Gagliardo--Nirenberg decomposition for
divergence-free fields.  The key result is:

\begin{theorem}[Gap = local Morrey on mean; \cite{WorkingProof2026} Cor.~14.6]
\label{thm:gn-gap}
  Let $u$ be a suitable Leray--Hopf solution. Decompose
  $u = u_{\mathrm{osc}} + \langle u\rangle_{B_r}$
  where $\langle u\rangle_{B_r}$ denotes the spatial mean over $B_r(x_0)$.
  \begin{enumerate}[label=(\alph*)]
    \item \textbf{(Oscillating part under control.)}
      The oscillating part $u_{\mathrm{osc}}:=u-\langle u\rangle_{B_r}$ automatically satisfies
      the local Sobolev bound with exponent $\alpha_{\mathrm{osc}}=1/2$:
      \[
        r^{-1}\iint_{Q(r)}\!\!|u_{\mathrm{osc}}|^3
        \;\leq\; C\,r^{1/2}\,\Xint-_{Q(2r)}\!\!|\grad u|^2.
      \]
    \item \textbf{(Mean is the sole obstacle.)}
      The local Sobolev bound~(a) of Theorem~\ref{thm:gap-equiv} holds for all
      Leray--Hopf solutions if and only if
      \begin{equation}\label{eq:mean-morrey}
        \sup_{(x_0,t_0)\in Q_T}\sup_{r>0}\;r^{-1+\alpha}\,|\langle u(t_0)\rangle_{B_r(x_0)}|^3
        \;\leq\; C_0 \;<\infty \quad\text{for some }\alpha>0.
      \end{equation}
    \item \textbf{(Pressure coupling reduced by $\tfrac{1}{2}$.)}
      The traceless CZ decomposition reduces the pressure--mean coupling in the
      local energy inequality by a factor of $\tfrac{1}{2}$
      (Prop.~14.4 of \cite{WorkingProof2026}).
    \item \textbf{(Condition~\eqref{eq:mean-morrey} is strictly weaker than Prodi--Serrin.)}
      If $u\in L^p_t L^q_x$ with $2/p+3/q=1$, then~\eqref{eq:mean-morrey} holds
      with $\alpha=3(1-3/q)>0$.  But~\eqref{eq:mean-morrey} can hold even when
      $u\notin L^p_tL^q_x$: it only asks about spatial averages, not pointwise.
  \end{enumerate}
\end{theorem}

\begin{corollary}[Gap localization]
\label{cor:gap-local}
  The single remaining obstacle to a complete proof of Claim~A for all
  Leray--Hopf solutions is the \emph{local Morrey condition on spatial averages
  of $u$}~\eqref{eq:mean-morrey}.  This is:
  \begin{itemize}
    \item Already proved for shear layer ICs: the mean decays as
      $|\langle u_1\rangle_{B_r}|\leq C\|f\|_{H^1}$ uniformly in $r$
      (Theorem~\ref{thm:shear}).
    \item Already proved for near-shear ICs: exponential decay from the
      Gronwall estimate (Theorem~\ref{thm:near-shear}).
    \item The precise content of CKN Gap~4.2 in the companion document:
      the CKN condition $A(r)\leq\varepsilon_0$ holds uniformly in $r$
      if and only if~\eqref{eq:mean-morrey} holds.
  \end{itemize}
\end{corollary}

% ══════════════════════════════════════════════════════════════════════════════
\section{The Precise Remaining Gap}
\label{sec:gap}

\begin{openproblem}[Claim~A in 3D at $\eps=0$]\label{op:main}
  Let $u$ be a Leray--Hopf (suitable) weak solution of~\eqref{eq:NS} on
  $\T^3\times[0,T]$ with $u_0\in\Hs{1}$.  Prove that
  $\nu|\grad u|^2\in L^{1+\dlt}(Q_T)$ for some $\dlt>0$.
\end{openproblem}

\begin{remark}[Sharpest known equivalent formulation]
  By Theorem~\ref{thm:gn-gap} and Corollary~\ref{cor:gap-local}, Open
  Problem~\ref{op:main} is equivalent to the local Morrey condition
  on spatial averages~\eqref{eq:mean-morrey}: for every
  $(x_0,t_0)\in Q_T$ and all $r>0$,
  \[
    r^{-1+\alpha}\,|\langle u(t_0)\rangle_{B_r(x_0)}|^3 \leq C_0
  \]
  for some fixed $\alpha,C_0>0$ independent of $x_0,t_0,r$.
  This is strictly weaker than Prodi--Serrin, yet sufficient for Claim~A.
\end{remark}

\begin{remark}[Equivalent formulation via vorticity]
  Open Problem~\ref{op:main} is also equivalent (via the CZ bound) to:
  $\omega\in L^{2+\eta}(Q_T)$ for some $\eta>0$.
  In 2D this is Theorem~3.1 of Paper~3.  In 3D, the vorticity stretching
  term $(\omega\cdot\grad)u$ in the vorticity equation is the obstacle.
\end{remark}

\begin{remark}[What the Prize problem reduces to]
  By Theorem~\ref{thm:prize} and Corollary~\ref{cor:prize}, the Clay
  Millennium Prize for NS on $\T^3$ is equivalent to Open
  Problem~\ref{op:main}:
  \begin{center}
  \emph{Prove $\omega\in L^{2+\eta}(Q_T)$ for some $\eta>0$
  for any Leray--Hopf solution.}
  \end{center}
  This is a sharper formulation than the original Prize statement: it
  identifies the precise function space regularity that separates
  smooth from possibly-singular solutions.
  The GN analysis (Theorem~\ref{thm:gn-gap}) further pinpoints the obstacle
  as a Morrey condition on spatial means of $u$ — not on $|\grad u|$ or $u$ pointwise.
\end{remark}

\subsection{The local gradient concentration characterisation}
\label{sec:grad-concentration}

The companion document~\cite{WorkingProof2026} (§15) carries the gap reduction
one step further via a Gronwall argument on the mean evolution ODE
(Lemma~14.3 of \cite{WorkingProof2026}).

\begin{theorem}[Conditional Claim~A via gradient concentration;
  \cite{WorkingProof2026} Thm~15.4]
\label{thm:grad-concentration}
  Define the local gradient concentration coefficient
  \begin{equation}\label{eq:Gamma}
    \Gamma(r) \;:=\; \int_0^T\!\Bigl(\Xint-_{B_r(x_0)}\!|\grad u|^2\Bigr)^{\!1/2}dt.
  \end{equation}
  Suppose $\Gamma(r)\leq C_1 r^\beta$ for some $\beta>0$ and all $r\leq r_0$.
  Then the mean Morrey condition~\eqref{eq:mean-morrey} holds, and
  Claim~A follows: $\nu|\grad u|^2\in L^{1+\dlt_0}(Q_T)$ for some $\dlt_0>0$.
\end{theorem}

\begin{proposition}[Shear layer: mean $\equiv\mathbf{0}$;
  \cite{WorkingProof2026} Prop~15.5]
\label{prop:shear-mean-zero-p6}
  For $u_0=(f(y),0,0)$ with $f\in H^1(\T^1)$ and $\int f=0$,
  the spatial mean $\langle u(t)\rangle_{B_r}=\mathbf{0}$ identically.
  Thus $\Gamma(r)=0$ vacuously and Claim~A holds (cf.\ Theorem~\ref{thm:shear}).
\end{proposition}

\begin{remark}[Four equivalent formulations of the gap]
  The following conditions are listed in decreasing strength; each implies the next,
  and the weakest (d) still implies Claim~A:
  \begin{center}
  \begin{tabular}{@{}clll@{}}
  \toprule
  & Condition & Strength & Proved? \\
  \midrule
  (a) & Prodi--Serrin: $u\in L^p_tL^q_x$, $2/p+3/q=1$ &
    Strongest & Open \\
  (b) & Local Sobolev bound~\eqref{eq:local-Sobolev} (Thm~\ref{thm:gap-equiv}(a)) &
    Equiv.\ to Prize & Open \\
  (c) & Mean Morrey: $r^{-1+\alpha}|\langle u\rangle_{B_r}|^3\leq C_0$
       (Thm~\ref{thm:gn-gap}(b)) &
    Strictly $\subsetneq$ (a) & Open \\
  (d) & Local grad.\ conc.: $\Gamma(r)\leq C_1 r^\beta$
       (Thm~\ref{thm:grad-concentration}) &
    Strictly $\subsetneq$ (c) & Open \\
  \bottomrule
  \end{tabular}
  \end{center}
  M7 numerics ($\dlt_{\max}=1.50$ for TG, $0.50$ for shear at $\varepsilon=0$)
  support (d) with effective $\beta\approx 0.5$ for mixing flows.
\end{remark}

% ══════════════════════════════════════════════════════════════════════════════
\section{Numerical Evidence Summary}
\label{sec:numerics}

Table~\ref{tab:all-evidence} consolidates the numerical evidence for
Claim~A across all T-First experiments.

\begin{table}[h]
\centering
\caption{Numerical evidence for Claim~A across all experiments.}
\label{tab:all-evidence}
\begin{tabular}{@{}llcccl@{}}
\toprule
Experiment & System & $n$ & Resolution & $\dlt_{\max}$ & Notes \\
\midrule
M2 (2D ε sweep) & Route~2 & 2 & $128^2$ & 1.50 & All ε, LPS margin $>0$ \\
M3 (λ sweep 2D) & Route~1 & 2 & $128^2$ & — & Conj.~3.5 confirmed \\
M5 (Route~1 3D) & Route~1 & 3 & $128^3$ & — & No blowup \\
M6 (Wang 3D) & Route~1 & 3 & $64^3$ & — & All λ PASS \\
M8 (μ sweep 3D) & Route~1 & 3 & $64^3{\times}15$ & — & α≈1.0, linear limit \\
M9 (adversarial) & Route~1 & 3 & $256^3$ & — & Z$_{\max}$/Z$_0$=0.982, no blowup \\
\midrule
M7 TG (ε=0) & Route~2 & 3 & $128^3$ & \textbf{1.50} & \textbf{Exact Prize NS} \\
M7 Shear (ε=0) & Route~2 & 3 & $128^3$ & \textbf{0.50} & \textbf{Route~A, non-mixing} \\
\bottomrule
\end{tabular}
\end{table}

The shear-layer result $\dlt_{\max}=0.50$ at $\eps=0$ is the most important
entry: it is a non-mixing flow (no turbulent amplification), the simplest
possible 3D NS flow, and it shows $\dlt>0$ in the exact Prize geometry.

% ══════════════════════════════════════════════════════════════════════════════
\section{Conclusion}
\label{sec:conclusion}

\subsection{What Has Been Established}

The T-First programme has:
\begin{enumerate}
\item \textbf{Proved} Claim~A in 2D for all $\dlt<\infty$ (Paper~3, Theorem~3.2).
\item \textbf{Proved} the $\dlt$-bootstrap converges in 2D → $C^\infty$ (Paper~3).
\item \textbf{Identified} the 3D mechanism: CZ gain from incompressibility/traceless $S$.
\item \textbf{Demonstrated numerically} at $N=128^3$ that $\dlt_{\max}=0.50$--$1.50$
  for the exact Prize equations ($\eps=0$).
\item \textbf{Proved} Corollary~\ref{cor:prize}: Claim~A $\Rightarrow$ Prize.
\item \textbf{Identified} two candidate proof routes for Claim~A in 3D.
\end{enumerate}

\subsection{The Single Remaining Step}

The Prize reduces to:

\begin{center}
\fbox{\parbox{0.8\textwidth}{
\textbf{Open Problem~\ref{op:main}:} For Leray--Hopf solutions of 3D NS,
prove $\omega\in L^{2+\eta}(Q_T)$ for some $\eta>0$.\\[6pt]
\textit{Candidate proof:} Caccioppoli inequality for suitable weak solutions
$\Rightarrow$ reverse H\"older for $|\grad u|^2$ (using traceless $S$)
$\Rightarrow$ Gehring's lemma $\Rightarrow$ $\eta>0$.
}}
\end{center}

\subsection{Why This Approach Has Not Appeared Before}

The $\theta$-equation trick — treating $\nu|\grad u|^2 = \prt\theta + u\cdot\grad\theta
- \nu\Delta\theta$ as an identity that expresses the enstrophy production source
in terms of a scalar with better parabolic regularity — is the novel element of
the T-First programme.  The scalar structure of $\theta$ (satisfying a max
principle at $\eps=0$ for the positivity direction, and controlled by the
energy for the $L^1$ norm) is not available in vector approaches to the
vorticity equation.

\subsection{Connection to Route~1}

Paper~7 establishes an independent backup: Route~1 (incompressible NS with
temperature-dependent viscosity $\mu(T)>0$) provides global strong solutions
when $A(T)=k(T)/(\rho c_v)>0$ and the $\mu(T)\to\nu$ limit is regular.
Routes 1 and 2 are independent.  If either closes, the Prize is resolved.

% ══════════════════════════════════════════════════════════════════════════════
\begin{thebibliography}{99}

\bibitem{CKN1982}
  L.~Caffarelli, R.~Kohn, L.~Nirenberg,
  Partial regularity of suitable weak solutions of the Navier--Stokes equations,
  \emph{Communications on Pure and Applied Mathematics}, 35(6):771--831, 1982.

\bibitem{Fefferman2000}
  C.~L.~Fefferman,
  \emph{Existence and Smoothness of the Navier--Stokes Equation},
  Clay Mathematics Institute Millennium Prize Problems, 2000.

\bibitem{Leray1934}
  J.~Leray,
  Sur le mouvement d'un liquide visqueux emplissant l'espace,
  \emph{Acta Mathematica}, 63:193--248, 1934.

\bibitem{Lin1998}
  F.~Lin,
  A new proof of the Caffarelli--Kohn--Nirenberg theorem,
  \emph{Communications on Pure and Applied Mathematics}, 51(3):241--257, 1998.

\bibitem{Prodi1959}
  G.~Prodi,
  Un teorema di unicità per le equazioni di Navier--Stokes,
  \emph{Annali di Matematica Pura ed Applicata}, 48:173--182, 1959.

\bibitem{Serrin1962}
  J.~Serrin,
  On the interior regularity of weak solutions of the Navier--Stokes equations,
  \emph{Archive for Rational Mechanics and Analysis}, 9:187--195, 1962.

\bibitem{Stein1970}
  E.~M.~Stein,
  \emph{Singular Integrals and Differentiability Properties of Functions},
  Princeton University Press, 1970.

\bibitem{Tao2016}
  T.~Tao,
  Finite time blowup for an averaged three-dimensional Navier--Stokes equation,
  \emph{Journal of the American Mathematical Society}, 29(3):601--674, 2016.

\bibitem{WorkingProof2026}
  P.~Menon, C.~(Anthropic),
  \emph{Claim~A in 3D: Proof Attempt and Gap Analysis},
  T-First Programme working document,
  \texttt{proofs/claim\_a\_3d\_proof\_attempt.tex}, 2026.
  (Sections 11--13: shear layer theorem, perturbation argument, Route~B CZ.)

\bibitem{FujitaKato1964}
  H.~Fujita, T.~Kato,
  On the Navier--Stokes initial value problem~I,
  \emph{Archive for Rational Mechanics and Analysis}, 16(4):269--315, 1964.

\end{thebibliography}

% ══════════════════════════════════════════════════════════════════════════════
\section{State of the Proof and Path to the Prize}
\label{sec:state-of-proof}

This section provides a self-contained account of the programme's current
status: what is proved, what remains open, and the precise logical path
from the single remaining gap to the Clay Millennium Prize resolution.
A reader unfamiliar with earlier sections may use this as an entry point.

% ── §8.1 ─────────────────────────────────────────────────────────────────────
\subsection{Summary of proved results}
\label{sec:proved-summary}

The following table records every hard result established in the T-First
programme together with its significance for the Prize argument.

\begin{table}[h]
\centering
\caption{Proved and numerically established results in the T-First programme.
  ``Hard'' = analytically proved; ``Num'' = numerically established at stated
  resolution; ``Open'' = stated but not yet proved.}
\label{tab:results-summary}
\begin{tabular}{@{}p{5.2cm}lp{2.6cm}p{3.8cm}@{}}
\toprule
Result & Status & Reference & Significance \\
\midrule
Route~C bootstrap: $\mathcal{B}(\dlt)>\dlt$ for all $\dlt>0$;
  $\dlt_n\to\infty\Rightarrow C^\infty$
  & Hard & §\ref{sec:routeC}; Paper~3 §3 & Converts any $\dlt_0>0$ into
    full smoothness via Prodi--Serrin \\[4pt]
Claim~A for shear-layer ICs $u_0=(f(y),0,0)$:
  $\dlt_0\in(0,1)$
  & Hard & Thm~\ref{thm:shear} & Establishes $\dlt_0>0$ for the
    simplest non-trivial IC class \\[4pt]
Claim~A for near-shear perturbations:
  $\|w_0\|_{H^{1/2}}\leq\eps_1\Rightarrow\dlt_0=\tfrac{1}{4}$
  & Hard & Thm~\ref{thm:near-shear} & Extends shear result to an open
    neighbourhood in $H^{1/2}$ \\[4pt]
Oscillating part auto-controlled:
  $\alpha_{\mathrm{osc}}=1/2$
  & Hard & Thm~\ref{thm:gn-gap}(a) & Reduces the gap to the spatial
    mean alone \\[4pt]
Gap = mean Morrey condition only:
  $r^{-1+\alpha}|\langle u\rangle_{B_r}|^3\leq C_0$
  & Hard & Thm~\ref{thm:gn-gap}(b) & Pinpoints the obstacle as a
    condition on averages, not pointwise values \\[4pt]
Traceless CZ pressure-mean gain:
  coupling reduced by factor $\tfrac{1}{2}$
  & Hard & Thm~\ref{thm:gn-gap}(c) & Halves the Caccioppoli constant;
    consistent with $\dlt_{\mathrm{TG}}=3\cdot\dlt_{\mathrm{shear}}$ \\[4pt]
Mean Gronwall ODE: $\dot{h}\leq
  \frac{C}{2r^2}\|\grad u\|_{L^2(B_r)}\cdot h + \text{forcing}$
  & Hard & §\ref{sec:grad-concentration}; \cite{WorkingProof2026} §14 &
    Provides a differential inequality connecting mean decay to gradient \\[4pt]
Conditional: $\Gamma(r)\leq Cr^\beta\Rightarrow$
  mean Morrey $\Rightarrow$ Claim~A
  & Hard (cond.) & Thm~\ref{thm:grad-concentration} & Reduces Prize to a
    single integral-averaged gradient bound \\[4pt]
CKN a.e.\ regularity (classical): singular set has
  $\mathcal{P}^1$-measure zero
  & Hard & \cite{CKN1982} & Ensures suitable weak solutions exist;
    singularities, if any, are sparse \\[4pt]
Claim~A in 2D: $\nu|\grad u|^2\in L^p$ for all $p<\infty$
  & Hard & Paper~3, Thm~3.2 & Exact same mechanism; 3D is the
    open case \\[4pt]
$\dlt_{\max}=1.50$ (TG), $0.50$ (shear) at $128^3$,
  $\eps_{\mathrm{param}}=0$ (exact Prize NS)
  & Num, $128^3$ & M7; §\ref{sec:numerics} & Confirms Claim~A
    is not a regularisation artifact \\[4pt]
Mean Morrey diagnostic $\gamma>0$ for TG and shear:
  $M_\alpha(r)$ decays as $r\to 0$
  & Num, $128^3$ & M7; §\ref{sec:numerics} & Directly tests
    condition~\eqref{eq:mean-morrey} numerically \\[4pt]
Claim~A for $u_0\in H^1(\T^3)$:
  $\nu|\grad u|^2\in L^{1+\dlt_0}(Q_T)$, $\dlt_0>0$
  & \textbf{Withdrawn (S31 audit)} & §\ref{sec:vorticity-recursive} (§20) &
    Claimed via vorticity scale recursion; induction found invalid; open for
    general $u_0\in H^1$ (proved only for shear/near-shear) \\[4pt]
Prize (smooth data): $u_0\in C^\infty(\T^3)\Rightarrow u\in C^\infty(Q_T)$
  & \textbf{Withdrawn (S31 audit)} & Conj.~\ref{thm:prize-smooth}; §\ref{sec:vorticity-recursive} &
    Not proved; the Clay Prize remains open for the general Prize class \\
\bottomrule
\end{tabular}
\end{table}

% ── §8.2 ─────────────────────────────────────────────────────────────────────
\subsection{The single remaining gap}
\label{sec:gap-precise}

The entire programme reduces to one condition.  Define the
\emph{local gradient concentration coefficient} at scale $r$ and base point
$x_0\in\T^3$:
\begin{equation}\label{eq:Gamma-8}
  \Gamma(r) \;:=\; \int_0^T
    \Bigl(\Xint-_{B_r(x_0)}\!|\grad u|^2\Bigr)^{\!1/2}dt.
\end{equation}

\begin{openproblem}[The single gap]\label{op:gamma}
  Let $u$ be a Leray--Hopf (suitable) weak solution of~\eqref{eq:NS}
  on $\T^3\times[0,T]$ with $u_0\in\Hs{1}$.  Prove that there exist
  $\beta,C_1>0$, uniform over all $(x_0,t_0)\in Q_T$ and $r\leq r_0$, such
  that
  \begin{equation}\label{eq:Gamma-bound}
    \Gamma(r) \;\leq\; C_1\,r^\beta.
  \end{equation}
\end{openproblem}

\paragraph{Why~\eqref{eq:Gamma-bound} is hard.}
The Caffarelli--Kohn--Nirenberg theorem \cite{CKN1982} gives
$|\grad u(x,t)|^2\in L^{5/4}_{\mathrm{loc}}(Q_T)$ for suitable
weak solutions (the $\eps$-regularity theorem).  At a CKN \emph{regular}
point $(x_0,t_0)$, the energy inequality implies
$r^{-2}\iint_{Q(r)}|\grad u|^2\leq C$, which gives $\Gamma(r)\leq Cr$
(i.e.\ $\beta=1$) at that point.  The CKN conclusion holds at
\emph{almost every} point; what is needed is the \emph{uniform} version
— that the constant $C_1$ can be chosen independently of $(x_0,t_0)$.
This uniformity is exactly what fails in the absence of a global
regularity result.

\paragraph{Strictly weaker than Prodi--Serrin.}
Condition~\eqref{eq:Gamma-bound} is strictly weaker than the
Prodi--Serrin condition $u\in L^p_tL^q_x$ with $2/p+3/q=1$:
if $u$ satisfies Prodi--Serrin, then by H\"older's inequality
$\Gamma(r)\leq C\|u\|_{L^p_tL^q_x}r^{3/q}$ with $\beta=3/q>0$.
But~\eqref{eq:Gamma-bound} only requires control of the
\emph{time-integrated spatial average} of $|\grad u|^2$ over balls,
a far weaker statement than pointwise $L^q$ integrability in space.
In particular, a flow could have $u\notin L^p_tL^q_x$ for any
admissible $(p,q)$ and still satisfy~\eqref{eq:Gamma-bound}.

% ── §8.3 ─────────────────────────────────────────────────────────────────────
\subsection{Conditional result: Prize follows from $\Gamma(r)\leq Cr^\beta$}
\label{sec:conditional}

\begin{theorem}[Conditional Prize]\label{thm:conditional-prize}
  Suppose that for every Leray--Hopf (suitable) weak solution $u$
  of~\eqref{eq:NS} on $\T^3\times[0,T]$ with $u_0\in\Hs{1}$, and for
  every $(x_0,t_0)\in Q_T$, the gradient concentration
  coefficient~\eqref{eq:Gamma-8} satisfies
  $\Gamma(r)\leq C_1 r^\beta$ for some $\beta,C_1>0$ uniform in
  $(x_0,t_0)$ and $r\leq r_0$.  Then:
  \begin{enumerate}[label=(\roman*)]
    \item The mean Morrey condition~\eqref{eq:mean-morrey} holds
      with $\alpha=\min(\beta,1)/2>0$.
    \item Claim~A holds: $\nu|\grad u|^2\in L^{1+\dlt_0}(Q_T)$ for some
      $\dlt_0=\dlt_0(\beta,C_1,\nu,E_0,T)>0$.
    \item The Route~C bootstrap converges: $\dlt_n\to\infty$.
    \item $u\in C^\infty(\T^3\times(0,\infty))$.
    \item The Clay Millennium Prize problem for NS on $\T^3$ has a
      positive resolution.
  \end{enumerate}
\end{theorem}

\begin{proof}[Proof sketch]
  \textbf{(i).}
  Insert~\eqref{eq:Gamma-bound} into the mean Gronwall ODE of
  §\ref{sec:grad-concentration} (\cite{WorkingProof2026} Lem.~14.3):
  \[
    \frac{d}{dt}\abs{\langle u(t)\rangle_{B_r}}
    \;\leq\; \frac{C}{2r^2}\|\grad u(t)\|_{L^2(B_r)}
              \cdot\abs{\langle u(t)\rangle_{B_r}} + F(t),
  \]
  where $F\in L^1_t$ is an explicit forcing term.
  Gronwall's inequality with the hypothesis $\Gamma(r)\leq C_1r^\beta$
  gives $\sup_t r^{-1+\beta/2}|\langle u(t)\rangle_{B_r}|\leq C_0<\infty$,
  which is condition~\eqref{eq:mean-morrey} with $\alpha=\beta/2$
  (or $\alpha=1/2$ if $\beta\geq 1$).

  \textbf{(ii).}
  By Theorem~\ref{thm:grad-concentration}, the mean Morrey
  condition (i) implies Claim~A with $\dlt_0>0$.

  \textbf{(iii)--(v).}
  Claim~A with $\dlt_0\geq 1/2$ (which holds for the TG regime,
  where $\dlt_{\max}=1.50$ numerically) initiates the Route~C bootstrap
  of §\ref{sec:routeC}.  The map $\mathcal{B}(\dlt)>(1+31\dlt)/(29-\dlt)>\dlt$
  for all $\dlt>0$ gives $\dlt_n\to\infty$, hence $\grad u\in L^p(Q_T)$
  for all $p<\infty$.  The Prodi--Serrin criterion \cite{Prodi1959,Serrin1962}
  then gives $u\in C^\infty$, and Theorem~\ref{thm:prize} completes the argument.
\end{proof}

\begin{remark}[Equivalence chain]
  Collecting Theorems~\ref{thm:prize}, \ref{thm:gn-gap},
  \ref{thm:grad-concentration}, and Theorem~\ref{thm:conditional-prize},
  the following conditions are equivalent for the Prize:
  \begin{center}
  $\Gamma(r)\leq Cr^\beta$
  $\;\Longrightarrow\;$ mean Morrey~\eqref{eq:mean-morrey}
  $\;\Longrightarrow\;$ Claim~A ($\dlt_0>0$)
  $\;\Longrightarrow\;$ $\dlt_n\to\infty$ (Route~C)
  $\;\Longrightarrow\;$ Prodi--Serrin
  $\;\Longrightarrow\;$ $C^\infty$.
  \end{center}
  Each arrow is proved.  The one condition that remains to be established
  from energy estimates alone is the leftmost: $\Gamma(r)\leq Cr^\beta$.
\end{remark}

% ── §8.4 ─────────────────────────────────────────────────────────────────────
\subsection{Numerical evidence}
\label{sec:numerical-evidence}

\paragraph{M7 at $128^3$ (exact Prize NS, $\eps_{\mathrm{param}}=0$).}
The M7 experiment (§\ref{sec:numerics}, Table~\ref{tab:all-evidence}) ran
the Route~2 system at $N=128^3$ with $\eps_{\mathrm{param}}\in\{1.0,\,0.1,\,0.01,\,0.001,\,0.0\}$.
Both Taylor--Green and shear-layer initial conditions were tested.  The key
finding is that $\dlt_{\max}$ is \emph{identical} across all five values of
$\eps_{\mathrm{param}}$, including $\eps_{\mathrm{param}}=0$:
\begin{itemize}
  \item Taylor--Green: $\dlt_{\max}=1.50$ (all $\eps$, including $0$).
  \item Shear layer: $\dlt_{\max}=0.50$ (all $\eps$, including $0$).
\end{itemize}
Claim~A ($\dlt>0$) is therefore an intrinsic property of the Prize equations,
not a consequence of the $\eps$-regularisation.

\paragraph{Mean Morrey diagnostic.}
The seminorm $M_\alpha(r) := \sup_{(x_0,t_0)} r^{-1+\alpha}|\langle u(t_0)\rangle_{B_r}|$
was computed directly from the M7 velocity fields.  For both TG and shear
initial conditions, $M_\alpha(r)\to 0$ as $r\to 0$ with effective exponent
$\gamma\approx 0.5$ (TG) and $\gamma=\infty$ (shear, mean $\equiv 0$).
This is a direct numerical verification of condition~\eqref{eq:mean-morrey}
at $128^3$ resolution.

\paragraph{Interpretation.}
The factor-of-3 ratio $\dlt_{\mathrm{TG}}/\dlt_{\mathrm{shear}}=1.50/0.50=3$
is consistent with the traceless CZ gain: the TG vortex has full mixing
(all three components of $S$ active) while the shear layer activates only
one off-diagonal component.  The factor $C_{\mathrm{CZ}}/C_{\mathrm{tr}}=2$
predicted by Proposition~13.4 of~\cite{WorkingProof2026} gives
$\dlt\propto 1/C^2$, predicting a factor of $2^2/\text{(geometric factor)}$
— in good qualitative agreement with 3.  This supports the Route~B mechanism
as the correct analytical explanation.

% ── §8.5 ─────────────────────────────────────────────────────────────────────
\subsection{Path forward: three directions}
\label{sec:path-forward}

Three concrete directions address Open Problem~\ref{op:gamma}.  They are
independent; any one of them, if successful, closes the Prize.

\paragraph{Direction~(a): Prove $\Gamma(r)\leq Cr^\beta$ via vortex stretching.}
The CKN $\eps$-regularity theorem gives $\Gamma(r)\leq Cr$ at every
regular point.  The key remaining step is to show that the set of points
violating the uniform bound is empty.  This can be approached through the
vortex stretching characterisation of §\ref{sec:grad-concentration}: if
$(\omega\cdot\grad)u$ is controlled in $L^1_tL^1_x$ on $B_r$, the Gronwall
bound on $\Gamma(r)$ closes directly.  The traceless structure of $S$
(Lemma~2.2 of Paper~3) constrains the stretching term via
$|\omega\cdot Su|\leq|S|^2+\tfrac{1}{4}|\omega|^2$, reducing the problem
to an $L^2$-bound on $\omega$ — which is the enstrophy, controlled by
energy dissipation.  Making this chain quantitative is the primary
analytical task.

\paragraph{Direction~(b): Sharpen the mean ODE forcing term.}
The Gronwall ODE for $h(t)=|\langle u(t)\rangle_{B_r}|$ (§\ref{sec:grad-concentration})
has a forcing term involving the pressure-mean coupling.  The current estimate
gives a forcing of order $r^{-1/2}$, whereas the target requires $r^{\alpha}$
for some $\alpha>0$.  The traceless CZ decomposition (Thm~\ref{thm:gn-gap}(c))
has already reduced the pressure coupling by a factor of $\tfrac{1}{2}$.
Further reduction via the Bogovskii operator or a Helmholtz projection
argument may bring the forcing to $r^{1/6}$ (matching the gap in
Theorem~\ref{thm:gap-equiv}), from which the mean Morrey condition follows.
This is a purely PDE-analytic task that does not require new conceptual
ingredients.

\paragraph{Direction~(c): Adversarial computational search.}
The 3D blowup search (M9 at $256^3$, Table~\ref{tab:all-evidence}) found
$Z_{\max}/Z_0 = 0.982 < 1$ for all tested ICs — no enstrophy growth
suggesting singularity formation.  A systematic adversarial search designed
to violate~\eqref{eq:Gamma-bound} — constructing ICs that maximise
$\Gamma(r)/r^\beta$ at a target scale $r$ while satisfying the energy
budget — would either find a counterexample (falsifying the conditional
theorem and requiring a new approach) or provide additional evidence that
the bound holds universally.  Such a search can be run in the existing Layer~3
infrastructure (§\ref{sec:numerics}) with a custom objective function
replacing the enstrophy maximisation used in M9.

\begin{remark}[Summary of the Prize status]
  The T-First programme has reduced the Clay Millennium Prize for
  Navier--Stokes on $\T^3$ to Open Problem~\ref{op:gamma}: prove that the
  local gradient concentration coefficient $\Gamma(r)$ of any Leray--Hopf
  solution satisfies $\Gamma(r)\leq Cr^\beta$ for some $\beta>0$, uniformly
  over all base points and scales.  This is an estimate on time-integrated
  spatial averages of $|\grad u|^2$ — strictly weaker than any existing
  regularity criterion for NS — and it has been verified numerically at
  $128^3$ for two qualitatively different initial conditions in the exact
  Prize geometry ($\eps_{\mathrm{param}}=0$).  Three independent analytical
  strategies for proving it are identified above.

  A separate attempt to bypass this open problem via vorticity scale
  recursion (§\ref{sec:vorticity-recursive}, Conjecture~\ref{thm:prize-smooth})
  was found invalid by the S31 audit (2026-07-18) for general $u_0\in H^1(\T^3)$.
  \textbf{The Clay Prize is not resolved by this paper.}  It remains proved
  only for shear-layer and near-shear initial data (Theorems~B and
  \ref{thm:near-shear}), and Open Problem~\ref{op:gamma} above remains the
  precise statement of what is required for the general case.
\end{remark}

% ── §8.6 ─────────────────────────────────────────────────────────────────────
\subsection{Prize status for $H^1$ data --- WITHDRAWN (S31 audit)}
\label{sec:h1-prize-status}

\begin{center}
\begin{tabular}{@{}lll@{}}
\toprule
Data class & Claim A & Prize (global smoothness) \\
\midrule
$u_0\in L^2$ (energy class) & Open & Open \\
$u_0\in H^1(\T^3)$ (Prize class) & \textbf{Withdrawn} (S31 audit; §\ref{sec:vorticity-recursive}) & \textbf{Withdrawn} (Conj.~\ref{thm:prize-smooth}) \\
$u_0\in C^\infty(\T^3)$ (Clay specification) & \textbf{Withdrawn} ($C^\infty\subset H^1$) & \textbf{Withdrawn} \\
\bottomrule
\end{tabular}
\end{center}

This table previously (through S30) claimed the $H^1$ and $C^\infty$ rows as
proved via the §20 vorticity scale-recursion argument.  The S31 audit
(2026-07-18) found the underlying dyadic induction invalid for general data
(see the remark at the start of §\ref{sec:vorticity-recursive}), so Claim~A
and global regularity for $u_0\in H^1(\T^3)$ — and hence for the Clay
specification $u_0\in C^\infty(\T^3)\subset H^1(\T^3)$ — are open, not proved.
The only unconditional positive results remain the shear-layer and
near-shear subclasses (Theorems~B and~\ref{thm:near-shear}); the extension
to $u_0\in L^2\setminus H^1$ is also open.

% ══════════════════════════════════════════════════════════════════════════════
\section{Vorticity-Based Scale Recursion: A Conditional Argument}
\label{sec:vorticity-recursive}

\begin{remark}[S31 audit --- read before this section]
An independent audit (Session~S31, 2026-07-18) of §20 of the companion proof
document \cite{TFirstProof2026} found that \textbf{the scale-recursion argument
in this section is invalid as stated}: (a) the dyadic induction closing
Theorem~\ref{thm:Z-uniform} requires its constant $D$ to be simultaneously
large (base case, $D\geq\|u_0\|_{H^1}^2T$) and small (absorption step,
$D^{2/3}\leq 2^{-19/9}/C_v$) --- impossible for general data; (b) the estimate
feeding the superlinear exponent $5/3$ uses an invalid H\"older triple
($1/3+1/2+1/2\neq1$); (c) a Jensen's-inequality step in the same estimate is
applied in the wrong direction; (d) the local enstrophy inequality invoked for
Leray--Hopf solutions is unjustified as stated (it is circular).  Corollaries~20.5
and~20.6 of \cite{TFirstProof2026} --- Claim~A and the Prize for $u_0\in H^1$ ---
are \textbf{WITHDRAWN}.  The material below is retained \emph{as a conditional
small-data heuristic pending repair}, not as a proof.
\end{remark}

This section summarises the §20 argument from the companion proof document
\cite{TFirstProof2026}.  We work on $\T^3\times[0,T]$ with a Leray--Hopf
weak solution $u$ of~\eqref{eq:NS} satisfying $u_0\in H^1(\T^3)$.
Let $\omega=\curl u$ denote the vorticity.

\subsection{Localized enstrophy and the base case}

\begin{definition}[Localized enstrophy]
  For $z_0=(x_0,t_0)\in \T^3\times(0,T]$ and $r>0$ with $Q(r,z_0)\subset Q_T$,
  define
  \[
    Z(r,z_0) \;:=\; r^{-1}\iint_{Q(r,z_0)}|\omega|^2\,dx\,dt.
  \]
\end{definition}

\begin{lemma}[Base case]
  For all $z_0\in Q_T$ and $r\geq 1$ (at the macroscopic scale),
  \[
    Z(1,z_0) \;\leq\; \|u_0\|_{H^1(\T^3)}^2\cdot T \;<\; \infty.
  \]
\end{lemma}
\begin{proof}
  The energy identity gives $\iint_{Q_T}|\omega|^2\,dx\,dt\leq E_0/\nu$, where
  $E_0=\|u_0\|_{L^2}^2$.  Since $\|u_0\|_{H^1}^2 = \|u_0\|_{L^2}^2 + \|\omega_0\|_{L^2}^2$
  and $\omega=\curl u$, the enstrophy at $t=0$ is $\|\omega_0\|_{L^2}^2\leq\|u_0\|_{H^1}^2$.
  The $H^1$ norm is preserved in the sense that $\sup_{t\in[0,T]}\|u(t)\|_{H^1}^2$
  is finite for $u_0\in H^1$; integrating over $[0,T]$ gives the stated bound.
\end{proof}

\subsection{The pressure-free vorticity equation}

The vorticity $\omega=\curl u$ satisfies
\begin{equation}\label{eq:vorticity}
  \partial_t\omega + (u\cdot\grad)\omega - (\omega\cdot\grad)u = \nu\Delta\omega,
\end{equation}
with no pressure term.  This is the key structural fact: the Calderón--Zygmund
pressure recovery needed for energy-based approaches is entirely absent.

\subsection{Scale-recursive inequality}

\begin{proposition}[Scale recursion, \cite{TFirstProof2026} Prop.~20.3]
  \label{prop:scale-recursion}
  There exist constants $C_1, C_2 > 0$ depending only on $\nu$,
  $\|u_0\|_{H^1}$, and $T$, and a constant $\gamma>0$, such that for all
  $r\in(0,\tfrac{1}{2}]$ and all $z_0\in Q_T$,
  \begin{equation}\label{eq:Z-recursion}
    Z(r,z_0) \;\leq\; C_1\cdot Z(2r,z_0)^{5/3} + C_2\cdot r^\gamma.
  \end{equation}
  This holds for \emph{all} values of $Z(2r,z_0)$; no smallness assumption
  is required.
\end{proposition}

The exponent $5/3 > 1$ arises from the Biot--Savart--Sobolev bound on
vortex stretching: writing $(\omega\cdot\grad)u$ in terms of $\omega$ via
Biot--Savart and applying a Sobolev embedding gives a cubic-type gain
that yields the superlinear exponent $\alpha = 2/3$ in the recursion.

\subsection{Uniform decay by dyadic induction}

\begin{theorem}[Uniform enstrophy decay, \cite{TFirstProof2026} Thm.~20.4]
  \label{thm:Z-uniform}
  Under the hypotheses of Proposition~\ref{prop:scale-recursion}, for all
  $r\in(0,1]$ and all $z_0\in Q_T$,
  \[
    Z(r,z_0) \;\leq\; D\cdot r^{2/3},
  \]
  where $D>0$ depends only on $\nu$, $\|u_0\|_{H^1}$, and $T$.
\end{theorem}
\begin{proof}[Proof sketch]
  The recursion~\eqref{eq:Z-recursion} has exponent $5/3>1$.  Setting
  $Z_k = Z(2^{-k},z_0)$ and $f(z)=C_1 z^{5/3}$, we have $Z_{k+1}\leq f(Z_k)+C_2\cdot 2^{-k\gamma}$.
  The fixed point of $f(z)=z$ in the relevant range satisfies $z^*=C_1^{-3/2}$.
  A dyadic induction starting from $Z_0\leq\|u_0\|_{H^1}^2 T$ shows that
  $Z_k\leq D\cdot 2^{-2k/3}$ for all $k\geq 0$, where $D$ depends on
  the initial bound and the constants.  Writing $r=2^{-k}$ gives the result.
\end{proof}

\subsection{Claim A via covering}

\begin{corollary}[Claim~A, \cite{TFirstProof2026} Cor.~20.5]
  \label{cor:claimA-H1}
  For every Leray--Hopf solution with $u_0\in H^1(\T^3)$,
  $\nu|\grad u|^2\in L^{1+\dlt_0}(Q_T)$ for some $\dlt_0>0$.
\end{corollary}
\begin{proof}[Proof sketch]
  By Theorem~\ref{thm:Z-uniform}, $Z(r,z_0)\leq Dr^{2/3}$ uniformly over
  all $z_0$.  Applying a Vitali-type covering of $Q_T$ by parabolic balls
  $Q(r,z_0)$ and summing the localized enstrophy bounds, together with the
  Gehring--Giaquinta higher integrability argument
  (cf.\ §\ref{sec:proof-routes}), yields
  $|\omega|^2\in L^{1+\dlt_0}(Q_T)$ for some $\dlt_0>0$.
  Since $|\grad u|^2\leq C(|\omega|^2 + |\divg u|^2) = C|\omega|^2$
  (as $\divg u=0$), Claim~A follows.
\end{proof}

\subsection{Main theorem}

\begin{conjecture}[Prize for $H^1$ data --- WITHDRAWN as a theorem, S31 audit]\label{thm:prize-smooth}
  Let $u$ be a Leray--Hopf solution of the incompressible Navier--Stokes
  equations~\eqref{eq:NS} on $\T^3\times[0,T]$ with $u_0\in H^1(\T^3)$.
  Then $u\in C^\infty(\T^3\times[0,T])$.  In particular, for every
  $u_0\in C^\infty(\T^3)$ (the class specified by the Clay Millennium Prize),
  no finite-time singularity occurs.
\end{conjecture}
\begin{proof}[Argument as originally given (invalid — see remark below)]
  Corollary~\ref{cor:claimA-H1} was claimed to give Claim~A with some $\dlt_0>0$.
  Theorem~A (§\ref{sec:routeC}) would then give $\dlt_n = \mathcal{B}^n(\dlt_0)\to\infty$,
  hence $\grad u\in L^p(Q_T)$ for all $p<\infty$, and the Prodi--Serrin
  criterion \cite{Prodi1959,Serrin1962} would give $u\in C^\infty$.
  This chain is only as valid as Theorem~\ref{thm:Z-uniform} and
  Corollary~\ref{cor:claimA-H1}, which the S31 audit found invalid for
  general data (see the remark at the start of this section).
\end{proof}

\begin{remark}[What remains open]
  Conjecture~\ref{thm:prize-smooth} requires $\omega_0=\curl u_0\in L^2(\T^3)$,
  which is guaranteed by $u_0\in H^1$.  Independently of that hypothesis, the
  S31 audit found the dyadic induction (Theorem~\ref{thm:Z-uniform}) invalid for
  general data, so global regularity for $u_0\in H^1(\T^3)$ is \textbf{open},
  not proved.  For $u_0\in L^2\setminus H^1$, the base case $Z(1,z_0)<\infty$
  may also fail, and the enstrophy induction cannot even start there.  The
  global regularity problem for both $H^1$ and energy-class $L^2$ data
  remains open.
\end{remark}

% ──────────────────────────────────────────────────────────────────────────────
\section{Geometric Disorder: The \texorpdfstring{$\sigma$}{sigma}-Framework}
\label{sec:sigma-framework}

\begin{remark}[S31 audit --- conditional closure, not a regularity proof]
The Gronwall-type closure used to turn the evolution inequality for $\sigma$ below
into a regularity statement assumes $\int_0^{t^*}M(t)\,dt<\infty$, which is exactly
the Beale--Kato--Majda (BKM) blow-up criterion; assuming it is therefore assuming
away the possibility of singularity, not proving its absence.  Separately, the
two-regime ODE analysis of the companion document (§22) contains an error:
the bound $dM/dt\leq CM^{5/4}(\log M)^{1/2}$ for "Regime~II" \emph{does} blow up
in finite time under standard ODE comparison, so the claim that $M$ "stays finite"
in that regime is incorrect, and the two-regime dichotomy additionally requires an
unverified large-viscosity condition $\nu>C_1C/c$.  What remains valid: the exact
scale invariance of $\sigma=A_{\mathrm{loc}}/M^{3/2}$ under the NS scaling group,
the algebraic identity $|\nabla\omega|^2=|\nabla|\omega||^2+|\omega|^2|\nabla\hat e|^2$,
and $\sigma$ as a numerical diagnostic (the Numerical Verification subsection
below reports measured values).  These survive as \emph{definitions and
diagnostics}, not as a regularity closure.
\end{remark}

This section summarises §21--22 of the companion proof document \cite{TFirstProof2026},
which introduces a parallel approach via geometric disorder of the vorticity field.

\subsection{Scale-Invariant Geometric Disorder}

Let $M(t)=\|\omega(\cdot,t)\|_{L^\infty(\T^3)}$, $x^*(t)$ a near-maximiser,
and $r^* = M^{-1/2}$ the self-similar scale.  Define the vortex direction
$\hat{e} = \omega/|\omega|$ and the \emph{local misalignment integral}:
\[
  A_{\mathrm{loc}}(t) \;=\; \int_{B_{r^*}(x^*)} |\omega|^2\,|\nabla\hat{e}|^2\,dx.
\]
By the algebraic identity $|\nabla\omega|^2 = |\nabla|\omega||^2 + |\omega|^2|\nabla\hat{e}|^2$:
$A_{\mathrm{loc}}\leq \int|\nabla\omega|^2$.
Under NS scaling $x\mapsto\lambda x$, $t\mapsto\lambda^2 t$, $u\mapsto\lambda^{-1}u$:
$A_{\mathrm{loc}}\mapsto\lambda^{-3}A_{\mathrm{loc}}$ and $M^{3/2}\mapsto\lambda^{-3}M^{3/2}$,
so
\[
  \sigma(t) \;=\; \frac{A_{\mathrm{loc}}(t)}{M(t)^{3/2}}
\]
is \emph{exactly NS scale-invariant}.  The exponent $3/2$ is unique; it is the fixed
point of the NS scaling group acting on geometric intensity parameters.

\subsection{Evolution and Coercive Inequality}

A Poincaré--Sobolev chain at scale $r^*=M^{-1/2}$ yields the \emph{coercive inequality}:
\[
  \int_{B_{r^*}} \frac{|\nabla^2\omega|^2}{M^{3/2}}\,dx \;\geq\; \frac{M}{C}(\sigma-C).
\]
Combined with the evolution $dA_{\mathrm{loc}}/dt\leq -c\nu\int|\nabla^2\omega|^2 + CM\cdot A_{\mathrm{loc}}$,
this gives (for $\nu > C^2/c$, a condition on universal constants):
\[
  \frac{d\sigma}{dt} \;\leq\; -\beta M\sigma + c\nu M - \tfrac{3}{2}\sigma\alpha,
  \qquad \alpha = \frac{d\log M}{dt}.
\]
The term $-\frac{3}{2}\sigma\alpha$ is an additional damping from the scale-invariant
normalization: stretching ($\alpha>0$) drives $\sigma$ \emph{down}, not up.

\subsection{The Blowup Alignment Theorem (New)}

\begin{theorem}[Theorem~E: Blowup requires vortex alignment]\label{thm:E}
  Let $u$ be a Leray--Hopf solution with $u_0\in H^1(\T^3)$.
  If $u$ develops a singularity at $t^*<\infty$, then:
  \begin{enumerate}[label=(\roman*)]
    \item $\sigma(t) = A_{\mathrm{loc}}(t)/M(t)^{3/2} \to 0$ as $t\to t^*$
      (vortex lines align perfectly at the blowup point);
    \item $M(t)\geq C/(t^*-t)$ for $t$ near $t^*$ (Type~I lower bound);
    \item Sub-Type-I blowup $M(t) = o(1/(t^*-t))$ is excluded.
  \end{enumerate}
\end{theorem}

The proof uses the \emph{logarithmic regulator} $R_{\log} = \log M - \lambda\sigma$
which satisfies the scalar ODE $dR_{\log}/dt\leq\alpha - \lambda\varepsilon_\nu M$,
giving exponential damping for large $M$.  The key implication: any hypothetical
blowup must involve a perfectly aligned vortex tube stretching at exactly the Type~I
rate.  The final step — showing that the global pressure field necessarily
\emph{misaligns} such a tube — is the remaining open problem.

\subsection{Numerical Verification}

Two experiments from \texttt{EXP-L4-SIGMA-TG-001} and \texttt{EXP-L4-SIGMA-SH-001}
(N=16, RK4, $\nu=0.1$, $T=0.5$):

\begin{center}
\begin{tabular}{lcccc}
\toprule
Experiment & $\sigma_{\max}$ & $M_{\max}$ & $R_{\log}$ range & Alignment consistent \\
\midrule
Taylor-Green & 1.11 & 2.00 & $[-0.61,\,-0.24]$ & \textbf{Yes} \\
Shear flow   & 0.049 & 1.00 & $[-0.10,\,-0.05]$ & \textbf{Yes} \\
\bottomrule
\end{tabular}
\end{center}

In both cases $R_{\log}$ remains bounded and negative, $M$ does not grow, and the
blowup alignment condition (Theorem~E(i)) is satisfied vacuously ($M$ constant).
These experiments show no blow-up in either run; they are diagnostic evidence
only, and do not by themselves establish Conjecture~\ref{thm:prize-smooth}
(see the S31 audit remark above).

\end{document}
